\section{Introduction}

\subsection{Scope}

This thesis makes one deliberate scope decision: it scores performance on a
single, well-defined task, not global surgical skill. Task scoring reduces
the problem to metrics that are objective and computable in real time --
completion time, path length, angular displacement, object state. Surgical
skill is broader, harder to define, and harder to measure objectively. Task
scoring is also solvable within a six-month thesis; general skill assessment
is not. This choice, not a limitation discovered later, shapes every design
decision in the chapters that follow.

\subsection{Goal}

The goal is real-time assessment of task performance on the da Vinci
Research Kit (dVRK), delivered as feedback during the task rather than
through post-hoc review. The target task is peg transfer: the trainee moves
a cylinder between pegs on a task pad under stereo endoscopic vision.

The system scores each trial automatically and objectively, without a human
rater in the loop. Four measures form the score: completion time, path
length of the instrument tip, angular displacement of the grasped object,
and object state (stationary, held, or dropped). Automatic scoring removes
inter-rater variability and makes the score available while the trial is
still running, not after the trainee has left the console.

The long-term vision extends past a single scored trial: a closed-loop
system that adapts task difficulty to the trainee's demonstrated
performance, without an instructor setting parameters by hand. The
adaptive-task design in Section~\ref{sec:conclusion} outlines one path
toward that vision; this thesis builds the real-time measurement layer
that any such closed loop would depend on.

\subsection{Motivation}

\paragraph{Real-time feedback accelerates learning.} Trainees who see their
performance while still performing the task can correct technique
immediately, instead of waiting for a debrief after the console session
ends. Immediate, task-relevant feedback is established in the motor-skill
learning literature as more effective for skill acquisition than delayed,
post-session review~\cite{TODO_motor_learning}. A real-time score turns
every trial into a feedback opportunity, not just a data point for later
analysis.

\paragraph{Automated capture catches what human observers miss.} An
expert observer watching a live trial cannot track every event at once --
a brief arm collision, incidental contact with the environment, a subtle
wrist over-rotation. These events are easy to miss in real time and hard
to reconstruct afterward from memory or video alone. Sensor-driven capture
logs every such event with precision, producing a complete record that
does not depend on what the observer happened to notice.

\paragraph{Multimodal data can surface metrics that redefine proficiency.}
Completion time is the traditional proxy for skill, but it collapses a lot
of information into one number. Combining kinematic, visual, and force data
opens the door to metrics that time alone cannot express -- motion
smoothness, instrument path efficiency, grasp stability. These metrics may
distinguish expert from novice behavior more precisely than time does, and
could inform how surgical proficiency is measured beyond this task.

\subsection{Contribution}

Prior work establishes the two ends of this problem separately: offline
kinematic analysis scores skill accurately but after the fact, and
real-time feedback systems exist but depend on a human expert to trigger
them (see Chapter~\ref{sec:background} for the detailed comparison). This
thesis closes that gap for one task: a multimodal, automated system that
scores dVRK peg transfer live, on the physical robot, with no human rater
and no post-hoc processing step in the feedback loop.
